\documentclass[11pt,a4paper]{article}

% ---------------------------------------------------------------------------
% Packages — deliberately conservative, so this compiles on a stock MiKTeX or
% TeX Live installation without downloading anything unusual.
% ---------------------------------------------------------------------------
\usepackage{lmodern}
\usepackage[T1]{fontenc}
\usepackage[utf8]{inputenc}
\usepackage[margin=2.4cm,headheight=14pt]{geometry}
\usepackage{xcolor}
\usepackage{amssymb}
\usepackage{enumitem}
\usepackage{booktabs}
\usepackage{longtable}
\usepackage{listings}
\usepackage{graphicx}   % \rotatebox, for the permission-matrix headings
\usepackage{fancyhdr}
\usepackage{titlesec}
\usepackage{parskip}
\usepackage{array}
\usepackage[hidelinks]{hyperref}

% ---------------------------------------------------------------------------
% Colours
% ---------------------------------------------------------------------------
\definecolor{ink}{HTML}{1E293B}
\definecolor{muted}{HTML}{64748B}
\definecolor{rule}{HTML}{CBD5E1}
\definecolor{shell}{HTML}{F1F5F9}
\definecolor{accent}{HTML}{DC2626}
\definecolor{good}{HTML}{047857}
\definecolor{warn}{HTML}{B45309}
\definecolor{warnbg}{HTML}{FEF3C7}
\definecolor{infobg}{HTML}{EFF6FF}
\definecolor{goodbg}{HTML}{ECFDF5}

\color{ink}

% ---------------------------------------------------------------------------
% Headers
% ---------------------------------------------------------------------------
\pagestyle{fancy}
\fancyhf{}
\renewcommand{\headrulewidth}{0.4pt}
\fancyhead[L]{\small\color{muted}IncidentFlow — Setup and Verification Manual}
\fancyhead[R]{\small\color{muted}\thepage}

\titleformat{\section}{\Large\bfseries\color{ink}}{\thesection.}{0.6em}{}
\titleformat{\subsection}{\large\bfseries\color{ink}}{\thesubsection}{0.6em}{}
\titlespacing*{\section}{0pt}{2.2ex plus 1ex minus .2ex}{1.2ex plus .2ex}

% ---------------------------------------------------------------------------
% Code listings
% ---------------------------------------------------------------------------
\lstdefinestyle{shellstyle}{
  basicstyle=\ttfamily\small,
  backgroundcolor=\color{shell},
  frame=single,
  rulecolor=\color{rule},
  framesep=6pt,
  xleftmargin=0pt,
  breaklines=true,
  breakatwhitespace=false,
  showstringspaces=false,
  columns=fullflexible,
  keepspaces=true,
}
\lstset{style=shellstyle}

% ---------------------------------------------------------------------------
% Callout boxes — built from primitives so no extra package is required.
% ---------------------------------------------------------------------------
\newcommand{\callout}[3]{%
  \par\medskip\noindent
  \fcolorbox{#1}{#2}{%
    \begin{minipage}{\dimexpr\linewidth-2\fboxsep-2\fboxrule\relax}
      \small #3
    \end{minipage}%
  }\par\medskip
}
\newcommand{\note}[1]{\callout{rule}{infobg}{#1}}
\newcommand{\warnbox}[1]{\callout{warn}{warnbg}{\textbf{Careful.} #1}}
\newcommand{\expect}[1]{\callout{good}{goodbg}{\textbf{You should see:} #1}}

% A checkbox item, for the verification lists.
\newcommand{\cbox}{\item[$\square$]}

% Inline shell command
\newcommand{\cmd}[1]{\texttt{\small #1}}

% ---------------------------------------------------------------------------
\begin{document}

\begin{titlepage}
\centering
\vspace*{3cm}

{\Huge\bfseries IncidentFlow}\\[0.4cm]
{\Large\color{muted} Setup and Verification Manual}\\[2cm]

{\large A step-by-step guide to running the platform\\
and checking every part of it yourself.}\\[1.5cm]

\rule{9cm}{0.4pt}\\[0.6cm]
{\color{muted}\small Written for someone who has never seen this project before.\\
No prior knowledge of Laravel, React or Docker is assumed.}\\[0.6cm]
\rule{9cm}{0.4pt}

\vfill
{\small\color{muted}Project directory: \texttt{D:\textbackslash IncidentFlow}}
\end{titlepage}

\tableofcontents
\newpage

% ===========================================================================
\section{Before you start}

\subsection{What this document is}

IncidentFlow is a platform for handling production incidents: staff report
problems, engineers coordinate the response, and managers track how reliable
the systems are over time.

This manual takes you from a cold machine to a running application, then walks
you through every screen and asks you to confirm what you see. Every section
ends with a checklist. If every box ticks, the system works.

\subsection{What is actually running}

Five separate programs cooperate. It helps to know what each one is for before
you start them.

\begin{center}
\begin{tabular}{@{}>{\raggedright\arraybackslash}p{3.2cm}>{\raggedright\arraybackslash}p{4.2cm}>{\raggedright\arraybackslash}p{6.6cm}@{}}
\toprule
\textbf{Program} & \textbf{Address} & \textbf{Its job} \\
\midrule
PostgreSQL & port 5432 & Stores everything. The single source of truth. \\
Redis & port 6379 & Carries live events and holds the job queue. \\
API (Laravel) & port 8000 & All the rules. Decides who may do what. \\
Realtime (Node) & port 3001 & Pushes live updates to open browsers. \\
Web app (React) & port 5173 & What you look at. \\
\bottomrule
\end{tabular}
\end{center}

\note{The web app on its own does nothing useful --- it asks the API for
everything. If a screen is empty, the question is almost always ``is the API
running?'' rather than ``is the page broken?''}

\subsection{Two ways to run it}

\begin{description}[leftmargin=1.6cm,style=nextline]
\item[Path A --- Docker] One command starts everything. Recommended if you
  have Docker Desktop installed. Skip to Section~\ref{sec:docker}.
\item[Path B --- Manual] Start each program yourself. More steps, but you see
  exactly what is happening and it needs no Docker. Section~\ref{sec:manual}.
\end{description}

To find out which you can use, open a terminal and type:

\begin{lstlisting}
docker --version
\end{lstlisting}

If that prints a version number, use Path A. If it says the command is not
recognised, use Path B.

\subsection{Opening a terminal}

Everything below is typed into a terminal.

\begin{itemize}[leftmargin=1.4cm]
\item \textbf{Windows}: press the Start button, type \cmd{powershell}, press Enter.
\item \textbf{macOS}: press \cmd{Cmd+Space}, type \cmd{terminal}, press Enter.
\end{itemize}

You will often need \emph{several} terminals open at once, because some
programs keep running and never give you the prompt back. Open a new one the
same way; that is normal and expected.

\newpage
% ===========================================================================
\section{Path A --- running it with Docker}
\label{sec:docker}

\subsection{Start everything}

\begin{lstlisting}
cd D:\IncidentFlow
docker compose up --build
\end{lstlisting}

The first run takes several minutes because it downloads and builds
everything. Later runs take seconds.

\expect{A long stream of coloured log lines that eventually settles down. Do
\emph{not} close this window --- closing it stops the application.}

\subsection{Confirm it came up}

Open a \emph{second} terminal and run:

\begin{lstlisting}
docker compose ps
\end{lstlisting}

\expect{A table listing \texttt{postgres}, \texttt{redis}, \texttt{mailpit}, \texttt{api},
\texttt{realtime}, \texttt{web}, \texttt{nginx}, \texttt{horizon} and
\texttt{scheduler}. The important column is \texttt{STATUS} --- everything
should say \texttt{running} or \texttt{healthy}.}

The \texttt{migrate} service is expected to say \texttt{exited (0)}. That is
correct: it creates the database tables and loads the demo data, then stops.
It is meant to run once, not forever.

\subsection{Open the application}

Go to \url{http://localhost:8080} in your browser.

Two other addresses are useful:

\begin{center}
\begin{tabular}{@{}ll@{}}
\toprule
\url{http://localhost:8025} & Mailpit --- every email the system sends lands here \\
\url{http://localhost:8080/horizon} & Horizon --- the background job dashboard \\
\bottomrule
\end{tabular}
\end{center}

\note{No real email is ever sent in this setup. Mailpit catches everything, so
you can read the notifications without anyone receiving them.}

Now jump to Section~\ref{sec:health}.

\newpage
% ===========================================================================
\section{Path B --- running it manually}
\label{sec:manual}

You will need four terminal windows. Keep them all open.

\subsection{Step 1 --- one-time preparation}

Do this once, ever. Skip it on later runs.

\begin{lstlisting}
cd D:\IncidentFlow\api
copy .env.local.example .env
composer install
php artisan key:generate
php artisan jwt:keys
php artisan migrate --seed
\end{lstlisting}

Taking those in turn:

\begin{description}[leftmargin=0.6cm,style=nextline,itemsep=2pt]
\item[\cmd{copy .env.local.example .env}] Creates the settings file. The
  \texttt{.env} file holds passwords and is deliberately never committed to
  version control.
\item[\cmd{composer install}] Downloads the PHP libraries the API needs.
\item[\cmd{php artisan key:generate}] Creates the encryption key for this
  installation.
\item[\cmd{php artisan jwt:keys}] Creates the pair of security keys used to
  sign login tokens. This prints a long line starting
  \texttt{JWT\_PUBLIC\_KEY=} --- you can ignore it for now.
\item[\cmd{php artisan migrate --seed}] Builds the database tables and fills
  them with ninety days of realistic demo incidents.
\end{description}

\expect{The last command prints a table of demo accounts and the line
\emph{``17 incidents across 90 days, 3 still open.''}}

Then prepare the web app, once:

\begin{lstlisting}
cd D:\IncidentFlow\web
npm install
\end{lstlisting}

\subsection{Step 2 --- start the API (terminal 1)}

\begin{lstlisting}
cd D:\IncidentFlow\api
php artisan serve --port=8000
\end{lstlisting}

\expect{\texttt{INFO Server running on [http://127.0.0.1:8000]}. Leave this
window open.}

\subsection{Step 3 --- start the web app (terminal 2)}

\begin{lstlisting}
cd D:\IncidentFlow\web
npm run dev
\end{lstlisting}

\expect{A short box showing \texttt{Local: http://localhost:5173/}. Leave this
window open.}

\subsection{Step 4 --- open it}

Go to \url{http://localhost:5173}.

\warnbox{At this point live updates are switched off. Everything else works
--- you just have to reload a page to see somebody else's changes. If that is
enough for you, skip to Section~\ref{sec:health}. To turn live updates on,
continue below.}

\subsection{Step 5 --- optional: live updates}

Live updates need Redis and the realtime service. Redis is already installed
at \texttt{D:\textbackslash redis}.

\textbf{Terminal 3 --- start Redis:}

\begin{lstlisting}
D:\redis\redis-server.exe D:\redis\incidentflow.conf
\end{lstlisting}

Check it answers, in any terminal:

\begin{lstlisting}
D:\redis\redis-cli.exe ping
\end{lstlisting}

\expect{\texttt{PONG}}

\textbf{Now tell the API to use it.} Open
\texttt{D:\textbackslash IncidentFlow\textbackslash api\textbackslash .env}
in a text editor and change these lines (add them if missing):

\begin{lstlisting}
REALTIME_ENABLED=true
REDIS_CLIENT=predis
REDIS_HOST=127.0.0.1
REDIS_PORT=6379
CACHE_STORE=redis
QUEUE_CONNECTION=redis
\end{lstlisting}

Stop the API (click its terminal, press \cmd{Ctrl+C}) and start it again.

\textbf{Terminal 4 --- start the realtime service.} This is one long command;
type it as a single line, or paste it.

\begin{lstlisting}
cd D:\IncidentFlow\realtime
$env:PORT="3001"
$env:REDIS_URL="redis://127.0.0.1:6379"
$env:JWT_PUBLIC_KEY_PATH="D:/IncidentFlow/api/storage/keys/jwt-public.pem"
$env:JWT_ISSUER="incidentflow-api"
$env:JWT_AUDIENCE="incidentflow-realtime"
$env:CORS_ORIGINS="http://localhost:5173"
npx tsx src/index.ts
\end{lstlisting}

\expect{Two lines of JSON, the second containing
\texttt{"message":"realtime service listening"}.}

\textbf{Terminal 5 --- optional, the job worker.} Without this, emails are
prepared but never sent.

\begin{lstlisting}
cd D:\IncidentFlow\api
php artisan queue:work redis
\end{lstlisting}

\newpage
% ===========================================================================
\section{Checking each part is alive}
\label{sec:health}

Before touching the interface, confirm each piece is responding. Run these in
a spare terminal.

\subsection{The API}

\begin{lstlisting}
curl http://localhost:8000/api/v1/health/live
\end{lstlisting}

\expect{\texttt{\{"status":"ok","service":"incidentflow-api",...\}}}

Now the deeper check, which tests the database and cache too:

\begin{lstlisting}
curl http://localhost:8000/api/v1/health/ready
\end{lstlisting}

\expect{\texttt{"status":"ready"} and \texttt{"degraded":false}.}

\note{If you skipped the Redis step, this will say \texttt{"degraded":true}.
That is not an error --- it means ``everything essential works, but live
updates are unavailable.'' The system is designed to keep running without
Redis rather than fall over.}

\subsection{The realtime service}

Only if you completed Step 5.

\begin{lstlisting}
curl http://localhost:3001/readyz
\end{lstlisting}

\expect{\texttt{\{"status":"ready","checks":\{"redis":"ok",...\}\}}}

\subsection{Verification checklist}

\begin{itemize}[leftmargin=1.2cm,itemsep=3pt]
\cbox The API answers \texttt{health/live} with \texttt{ok}.
\cbox The API answers \texttt{health/ready} with \texttt{ready}.
\cbox The web page loads and shows a sign-in form.
\cbox (If applicable) Redis answers \texttt{PONG}.
\cbox (If applicable) The realtime service answers \texttt{ready}.
\end{itemize}

\newpage
% ===========================================================================
\section{The demo accounts}

Six accounts exist. \textbf{The password for every one is} \cmd{incidentflow}.

They are not interchangeable --- each sees a genuinely different application.
That is the point of the next several sections.

\begin{center}
\begin{tabular}{@{}>{\raggedright\arraybackslash}p{6.1cm}>{\raggedright\arraybackslash}p{3cm}>{\raggedright\arraybackslash}p{4.9cm}@{}}
\toprule
\textbf{Email} & \textbf{Role} & \textbf{Can do} \\
\midrule
\texttt{admin@incidentflow.test} & Administrator & Everything, plus the audit log \\
\texttt{commander@incidentflow.test} & Incident Commander & Assign people, change severity, publish postmortems, export \\
\texttt{responder@incidentflow.test} & Responder & Move an incident through its stages, post updates \\
\texttt{responder2@incidentflow.test} & Responder & The same, again --- so there is somebody to assign \\
\texttt{reporter@incidentflow.test} & Reporter & Report incidents, comment \\
\texttt{viewer@incidentflow.test} & Viewer & Read only \\
\bottomrule
\end{tabular}
\end{center}

\note{Sign in and out often as you work through this manual. The quickest way
to understand the permission system is to look at the same screen as two
different people.}

\newpage
% ===========================================================================
\section{Guided tour}

Work through these in order. Each has instructions, then a checklist.

% ---------------------------------------------------------------------------
\subsection{The incident register}

Sign in as \texttt{reporter@incidentflow.test}.

You land on a table of incidents. Things to try:

\begin{enumerate}[leftmargin=1.2cm,itemsep=3pt]
\item Type \cmd{checkout} into the search box. The list narrows as you type.
\item Set \textbf{All severities} to \textbf{SEV1}. Only the most serious
  incidents remain.
\item Tick \textbf{Active only}. Anything already resolved disappears.
\item Clear the filters again.
\end{enumerate}

Look at the \textbf{Running / resolved in} column. For open incidents this is
a live clock showing how long the problem has been going on.

Look at the navigation bar. Count the tabs --- there should be \textbf{four}.

\begin{itemize}[leftmargin=1.2cm,itemsep=3pt]
\cbox 17 incidents are listed.
\cbox Searching narrows the list.
\cbox Filtering by severity works.
\cbox There are four navigation tabs, and \emph{no} Admin tab.
\cbox A \textbf{Report incident} button is visible.
\end{itemize}

\note{The missing Admin tab is deliberate. A Reporter has no permission to
read the audit log, so the tab does not exist for them --- it is not merely
greyed out.}

% ---------------------------------------------------------------------------
\subsection{Reporting an incident}

Still signed in as the Reporter, click \textbf{Report incident}.

Fill it in:

\begin{itemize}[leftmargin=1.2cm,itemsep=2pt]
\item \textbf{What is happening?} --- \emph{Card payments failing for UK customers}
\item \textbf{Severity} --- SEV-1
\item \textbf{Service} --- Checkout API
\item \textbf{What do we know?} --- anything you like
\end{itemize}

Read the amber warning before you submit. It tells you what choosing SEV-1
actually causes. Then click \textbf{Report incident}.

You are taken straight to the new incident.

\begin{itemize}[leftmargin=1.2cm,itemsep=3pt]
\cbox You are moved to the new incident's own page.
\cbox It has a reference like \texttt{INC-0018}.
\cbox A red \textbf{Open} badge and a \textbf{Postmortem required} badge appear.
\cbox The timeline already has one entry: \emph{``... reported the incident''}.
\cbox There are \textbf{no} buttons to change the status.
\end{itemize}

\warnbox{That last point is the important one. A Reporter cannot move an
incident forward. Try to find a way to do it --- you should not be able to.}

\subsubsection*{Check the emails went out}

If you are on Path A (Docker), open \url{http://localhost:8025}.

\expect{Four emails, one to each person who can be paged. There is
\emph{nothing} addressed to the Reporter or the Viewer --- neither can be
assigned to an incident, so neither is on the on-call rota.}

% ---------------------------------------------------------------------------
\subsection{Working the incident}

Sign out. Sign in as \texttt{responder@incidentflow.test}. Open the incident
you just created.

Notice you now have buttons. Count them --- there are \textbf{two}:
\textbf{Mark acknowledged} and \textbf{Mark resolved}.

\note{There is no ``Mark closed'' button. An incident cannot be closed before
it has been resolved, so that option is not offered. The buttons on screen are
generated from the rules, so you can never be shown an action that would be
refused.}

Click \textbf{Mark acknowledged}. Watch the buttons change to \textbf{Mark
mitigated} and \textbf{Mark resolved}.

Now scroll down and try the two writing areas:

\begin{itemize}[leftmargin=1.2cm,itemsep=2pt]
\item \textbf{Post an update} --- the public record of what is happening.
  Tick \emph{Safe for a customer status page} before posting one.
\item \textbf{Discussion} --- internal chat between responders.
\end{itemize}

\begin{itemize}[leftmargin=1.2cm,itemsep=3pt]
\cbox Exactly two status buttons appear before you click.
\cbox After acknowledging, the badge turns amber and reads \textbf{Acknowledged}.
\cbox The buttons change to Mitigated / Resolved.
\cbox The timeline gains \emph{``... acknowledged the incident''}.
\cbox A posted update shows a green \textbf{Public} tag when you tick the box.
\cbox Your comment appears in Discussion immediately.
\cbox The right-hand panel shows \textbf{Time to acknowledge} with a target.
\end{itemize}

% ---------------------------------------------------------------------------
\subsection{Commanding the incident}

Sign out. Sign in as \texttt{commander@incidentflow.test}. Open the same
incident and scroll to the right-hand column.

You now have two panels the Responder did not: \textbf{Responders} and
\textbf{Change severity}.

\subsubsection*{Assigning someone}

Open the \textbf{Assign someone...} dropdown and read the names.

\expect{Only people who can actually be paged. Ahmed Hassan (Reporter) and
Sofia Rossi (Viewer) are absent --- assigning a Viewer would look like
coverage while providing none.}

Pick someone. They appear in the list with a button to remove them again.

\subsubsection*{Changing severity}

This one is worth doing carefully, because the behaviour is deliberately
asymmetric.

\begin{enumerate}[leftmargin=1.2cm,itemsep=3pt]
\item Choose a \emph{higher} severity (e.g. SEV-1 to SEV-2 is lower; go the
  other way). No extra field appears --- making something more serious needs
  no justification.
\item Now choose a \emph{lower} severity, such as SEV-4. A \textbf{reason}
  box appears, and it is required.
\end{enumerate}

\note{Quietly downgrading an incident changes whether a postmortem is
required and how it appears in every later report. So a downgrade goes on the
record; an escalation does not need to.}

\begin{itemize}[leftmargin=1.2cm,itemsep=3pt]
\cbox The Responders and Change severity panels are visible.
\cbox The assignment dropdown excludes the Reporter and the Viewer.
\cbox An assigned person appears in the list with a remove button.
\cbox Raising severity asks for no reason.
\cbox Lowering severity demands a written reason.
\end{itemize}

% ---------------------------------------------------------------------------
\subsection{Live updates}
\label{sec:live}

\warnbox{This section only works if you completed Path A, or Path B Step 5.
Look at the top right of the screen: if it says a green \textbf{Live}, you are
ready. If it says amber \textbf{Reconnecting...}, skip this section.}

This is best seen with two browser windows side by side.

\begin{enumerate}[leftmargin=1.2cm,itemsep=3pt]
\item In your normal window, stay signed in as the Commander, on an incident.
\item Open a \textbf{private / incognito} window (\cmd{Ctrl+Shift+N} in
  Chrome). Go to the same address and sign in as
  \texttt{responder@incidentflow.test}. Open the \emph{same} incident.
\item Arrange the windows so you can see both.
\item In the \textbf{private} window, post a comment.
\item Watch the \textbf{other} window without touching it.
\end{enumerate}

\expect{The comment and a new timeline entry appear in the first window within
a second or two. You did not reload anything.}

\begin{itemize}[leftmargin=1.2cm,itemsep=3pt]
\cbox The header shows a green \textbf{Live} badge.
\cbox A change made in one window appears in the other unprompted.
\cbox The incident list also updates itself the same way.
\end{itemize}

% ---------------------------------------------------------------------------
\subsection{Resolving, and the postmortem}

As the Commander, click \textbf{Mark resolved}.

Notice the buttons become \textbf{Mark closed} and \textbf{Mark
acknowledged}. Resolved is not the end --- if the problem returns you reopen
it, and the system discards the old resolution time rather than flattering the
statistics.

Now scroll to the \textbf{Postmortem} section.

\begin{enumerate}[leftmargin=1.2cm,itemsep=3pt]
\item Read the four required sections. Each is labelled \emph{required ---
  blocks publishing} in amber, and each has a hint telling you what belongs
  there.
\item Fill in \textbf{Summary} only, then click \textbf{Start postmortem}.
\item Look at the \textbf{Publish} button --- it is greyed out, and a note
  below lists exactly what is still missing.
\item Now fill in Root cause, Impact and Resolution. Watch each label change
  from amber \emph{required} to green \emph{complete}.
\item Click \textbf{Save draft}, then \textbf{Publish}.
\end{enumerate}

\expect{A green banner naming who published it and when. Every box becomes
read-only, and the Save and Publish buttons disappear entirely.}

Now try to edit it. You cannot --- and neither can an administrator.

\begin{itemize}[leftmargin=1.2cm,itemsep=3pt]
\cbox Four sections are marked as required.
\cbox Publish is disabled while any of them is empty.
\cbox The message below names precisely which are missing.
\cbox Labels turn green as you complete each one.
\cbox Publish becomes available only when all four are filled.
\cbox After publishing, every field is locked and the buttons vanish.
\end{itemize}

% ---------------------------------------------------------------------------
\subsection{The metrics dashboard}

Click \textbf{Metrics}. Switch between the 7, 30 and 90 day buttons.

Read the two large numbers carefully. Under \textbf{Mean time to
acknowledge} there is a smaller line reading something like
\emph{p50 9m $\cdot$ p95 5h}.

\note{This is the most important idea on the page. ``p50'' is the middle
case: half of all incidents were picked up faster than that. ``p95'' is the
bad case: one in twenty took at least that long. An average on its own hides
both. If the average is 46 minutes but the middle case is 9 minutes, then
most incidents are handled quickly and a few are handled very slowly --- and
those are two completely different problems to fix.}

Scroll down for the SLA table and the two charts. Hover over a bar in
\textbf{Opened vs resolved} to see that day's figures.

\begin{itemize}[leftmargin=1.2cm,itemsep=3pt]
\cbox Four summary figures appear at the top.
\cbox Each average is accompanied by p50 and p95.
\cbox The SLA table shows a target and an attainment per severity.
\cbox Both charts draw, and the bar chart responds to hovering.
\cbox Changing the period changes the numbers.
\end{itemize}

% ---------------------------------------------------------------------------
\subsection{Services}

Click \textbf{Services}. These are the systems that can break.

\begin{itemize}[leftmargin=1.2cm,itemsep=3pt]
\cbox Five services are listed.
\cbox Each shows a tier and the team that owns it.
\cbox Services with live incidents show a red count.
\end{itemize}

% ---------------------------------------------------------------------------
\subsection{Exporting}

Still as the Commander, go back to \textbf{Incidents}. There is now an
\textbf{Export CSV} button.

Do this in order, because the point is what the file contains:

\begin{enumerate}[leftmargin=1.2cm,itemsep=3pt]
\item Without any filters, click \textbf{Export CSV}. Open the downloaded
  file in Excel. Count the rows.
\item Back in the browser, tick \textbf{Active only}. Note how many rows the
  table now shows.
\item Click \textbf{Export CSV} again and open that file.
\end{enumerate}

\expect{The second file contains exactly the rows that were on screen ---
not everything. The export always matches what you were looking at.}

\begin{itemize}[leftmargin=1.2cm,itemsep=3pt]
\cbox A file downloads, named after the organisation and the date.
\cbox It opens in Excel with readable column headings.
\cbox Filtering the screen changes what the file contains.
\end{itemize}

% ---------------------------------------------------------------------------
\subsection{Administration}

Sign out. Sign in as \texttt{admin@incidentflow.test}.

Count the navigation tabs --- there are now \textbf{five}. The \textbf{Admin}
tab has appeared.

Open it. The top half lists everyone and their role. Scroll down for the
audit log.

\expect{Every action you have taken during this tour, most recent first, with
who did it and when. Underneath each is a long identifier --- that is the
tracking number for that request, shared across all parts of the system.}

\begin{itemize}[leftmargin=1.2cm,itemsep=3pt]
\cbox Five navigation tabs, including Admin.
\cbox All six members are listed with their roles.
\cbox The audit log shows your actions from this session.
\cbox Each entry names a person.
\cbox Your CSV exports appear in the log too.
\end{itemize}

\note{Exports are recorded because taking customer data out of the system is
itself something worth being able to account for later.}

\newpage
% ===========================================================================
\section{Checking the promises yourself}

The sections above show the product working. These show it refusing to
misbehave. Each is a claim you can test.

% ---------------------------------------------------------------------------
\subsection{A Viewer really is read-only}

Sign in as \texttt{viewer@incidentflow.test}.

\begin{itemize}[leftmargin=1.2cm,itemsep=3pt]
\cbox No \textbf{Report incident} button anywhere.
\cbox Opening an incident offers no status buttons.
\cbox No comment box.
\cbox No Admin tab.
\end{itemize}

Now try to get around the interface. In the address bar, type
\url{http://localhost:5173/admin} directly.

\expect{You are not shown the audit log. The restriction lives in the server,
not in whether a button was drawn.}

% ---------------------------------------------------------------------------
\subsection{An incident cannot skip stages}

The interface never offers an illegal step, so this must be tested directly
against the API.

First get a login token (this is one long line):

\begin{lstlisting}
curl -X POST http://localhost:8000/api/v1/auth/login -H "Content-Type: application/json" -d "{\"email\":\"responder@incidentflow.test\",\"password\":\"incidentflow\"}"
\end{lstlisting}

Copy the long string after \texttt{"access\_token":}. Now try to close an
incident that is merely open --- replace \texttt{PASTE\_TOKEN} with it, and
\texttt{18} with any open incident's number:

\begin{lstlisting}
curl -X POST http://localhost:8000/api/v1/incidents/18/status -H "Authorization: Bearer PASTE_TOKEN" -H "X-Organization: northwind" -H "Content-Type: application/json" -d "{\"status\":\"closed\"}"
\end{lstlisting}

\expect{A refusal containing \texttt{"code":"incident.illegal\_transition"},
and helpfully a list of which steps \emph{are} allowed from here.}

% ---------------------------------------------------------------------------
\subsection{A double-submitted report does not create two incidents}

If someone taps ``Report'' twice on a bad connection, only one incident should
exist. Run this command \emph{twice}, unchanged:

\begin{lstlisting}
curl -X POST http://localhost:8000/api/v1/incidents -H "Authorization: Bearer PASTE_TOKEN" -H "X-Organization: northwind" -H "Content-Type: application/json" -H "Idempotency-Key: my-test-key-123" -d "{\"title\":\"Testing duplicate submission\",\"severity\":\"sev3\"}"
\end{lstlisting}

\expect{Both responses describe the \emph{same} incident, with the same
reference number. Check the incident list --- there is only one.}

\note{The \texttt{Idempotency-Key} is what makes this safe. The application
generates one automatically for every report you file through the interface.}

% ---------------------------------------------------------------------------
\subsection{History cannot be rewritten}

Look at any incident's timeline. There is no edit button and no delete button
--- not for a Responder, not for a Commander, not for an Administrator.

This is enforced in three separate places, so that removing one would not be
enough. The reason is simple: a postmortem is written from the timeline, so if
the timeline could be altered afterwards, nothing concluded from it could be
trusted.

\begin{itemize}[leftmargin=1.2cm,itemsep=3pt]
\cbox No timeline entry can be edited or deleted by anyone.
\cbox No audit log entry can be edited or deleted by anyone.
\cbox A published postmortem cannot be edited, even by its author.
\end{itemize}

% ---------------------------------------------------------------------------
\subsection{Repeated login attempts are throttled}

Run this six times quickly, with the wrong password on purpose:

\begin{lstlisting}
curl -X POST http://localhost:8000/api/v1/auth/login -H "Content-Type: application/json" -d "{\"email\":\"admin@incidentflow.test\",\"password\":\"wrong\"}"
\end{lstlisting}

\expect{The first few refuse politely. Around the sixth, you get
\texttt{"code":"rate\_limited"} and a \texttt{Retry-After} header.}

Also notice: guessing an address that does not exist gives \emph{exactly} the
same refusal as a wrong password. The login form cannot be used to discover
who has an account.

\newpage
% ===========================================================================
\section{Running the automated tests}

You do not need the application running for these. Each is independent.

\subsection{The API}

\begin{lstlisting}
cd D:\IncidentFlow\api
php vendor/bin/phpunit
\end{lstlisting}

\expect{\texttt{OK (104 tests, 518 assertions)} in green, after a few seconds.}

\subsection{The realtime service}

\begin{lstlisting}
cd D:\IncidentFlow\realtime
npm test
\end{lstlisting}

\expect{\texttt{Tests 26 passed (26)}}

\subsection{The web app}

\begin{lstlisting}
cd D:\IncidentFlow\web
npm test
\end{lstlisting}

\expect{\texttt{Tests 23 passed (23)}}

\subsection{The browser tests}

These drive a real browser through the whole application, so the system must
be running first.

\begin{lstlisting}
cd D:\IncidentFlow\web
npx playwright install chromium
npm run e2e
\end{lstlisting}

\expect{\texttt{7 passed}, taking a couple of minutes.}

\note{If you are on Path B rather than Docker, add
\cmd{E2E\_EXPECT\_TIMEOUT=45000} before the command. The development server
answers one request at a time, so pages load more slowly than the tests
expect by default.}

\begin{itemize}[leftmargin=1.2cm,itemsep=3pt]
\cbox API: 104 tests pass.
\cbox Realtime: 26 tests pass.
\cbox Web: 23 tests pass.
\cbox Browser: 7 tests pass.
\end{itemize}

\newpage
% ===========================================================================
\section{When something goes wrong}

\subsection{The page is blank or says it cannot load}

The API is almost certainly not running. Check its terminal window, then:

\begin{lstlisting}
curl http://localhost:8000/api/v1/health/live
\end{lstlisting}

If that fails, restart the API.

\subsection{``Address already in use''}

Something is already using that port --- often a copy of this application you
forgot to stop. On Windows:

\begin{lstlisting}
netstat -ano | findstr :8000
taskkill /F /PID <the number from the last column>
\end{lstlisting}

\subsection{The header says ``Reconnecting...'' forever}

Live updates are unavailable. In order, check:

\begin{enumerate}[leftmargin=1.2cm,itemsep=2pt]
\item Is Redis running? \cmd{D:\textbackslash redis\textbackslash redis-cli.exe ping}
\item Is the realtime service running? \cmd{curl http://localhost:3001/readyz}
\item Does the API's \texttt{.env} say \texttt{REALTIME\_ENABLED=true}?
\item Did you restart the API after changing that?
\end{enumerate}

\warnbox{This never stops the rest of the application working. Everything
except automatic updating continues normally --- you just reload to see
changes.}

\subsection{Every page says you are signed out}

The security keys are missing or were regenerated. Rebuild them:

\begin{lstlisting}
cd D:\IncidentFlow\api
php artisan jwt:keys --force
\end{lstlisting}

Then restart the API and the realtime service, and sign in again.

\subsection{Starting completely fresh}

This deletes all data and rebuilds the demo from scratch. It is safe --- there
is nothing real in here.

\begin{lstlisting}
cd D:\IncidentFlow\api
php artisan migrate:fresh --seed
\end{lstlisting}

\subsection{Stopping everything}

Docker: press \cmd{Ctrl+C} in its window, then \cmd{docker compose down}.

Manual: press \cmd{Ctrl+C} in each terminal window.

\newpage
% ===========================================================================
\appendix
\section{Command reference}

\begin{center}
\begin{tabular}{@{}>{\raggedright\arraybackslash}p{8.6cm}>{\raggedright\arraybackslash}p{7cm}@{}}
\toprule
\textbf{Command} & \textbf{What it does} \\
\midrule
\texttt{docker compose up -{}-build} & Start everything (Path A) \\
\texttt{docker compose down} & Stop everything \\
\texttt{docker compose logs -f api} & Watch the API's log \\
\midrule
\texttt{php artisan serve -{}-port=8000} & Start the API (Path B) \\
\texttt{php artisan queue:work redis} & Start the job worker \\
\texttt{php artisan migrate -{}-seed} & Create tables, load demo data \\
\texttt{php artisan migrate:fresh -{}-seed} & Wipe and rebuild everything \\
\texttt{php artisan jwt:keys -{}-force} & Regenerate the security keys \\
\texttt{php artisan incidentflow:prune -{}-dry-run} & Show what housekeeping would delete \\
\midrule
\texttt{npm run dev} & Start the web app \\
\texttt{npm test} & Run that project's tests \\
\texttt{npm run e2e} & Run the browser tests \\
\bottomrule
\end{tabular}
\end{center}

\section{Who can do what}

\begin{center}
\begin{tabular}{@{}>{\raggedright\arraybackslash}p{5.4cm}ccccc@{}}
\toprule
& \rotatebox{60}{Viewer} & \rotatebox{60}{Reporter} & \rotatebox{60}{Responder} & \rotatebox{60}{Commander} & \rotatebox{60}{Admin} \\
\midrule
View incidents        & $\checkmark$ & $\checkmark$ & $\checkmark$ & $\checkmark$ & $\checkmark$ \\
View metrics          & $\checkmark$ & $\checkmark$ & $\checkmark$ & $\checkmark$ & $\checkmark$ \\
Report an incident    &              & $\checkmark$ & $\checkmark$ & $\checkmark$ & $\checkmark$ \\
Comment               &              & $\checkmark$ & $\checkmark$ & $\checkmark$ & $\checkmark$ \\
Change status         &              &              & $\checkmark$ & $\checkmark$ & $\checkmark$ \\
Post updates          &              &              & $\checkmark$ & $\checkmark$ & $\checkmark$ \\
Assign responders     &              &              &              & $\checkmark$ & $\checkmark$ \\
Change severity       &              &              &              & $\checkmark$ & $\checkmark$ \\
Publish postmortems   &              &              &              & $\checkmark$ & $\checkmark$ \\
Export to CSV         &              &              &              & $\checkmark$ & $\checkmark$ \\
Manage services       &              &              &              &              & $\checkmark$ \\
Manage members        &              &              &              &              & $\checkmark$ \\
Read the audit log    &              &              &              &              & $\checkmark$ \\
Delete incidents      &              &              &              &              & $\checkmark$ \\
\bottomrule
\end{tabular}
\end{center}

Each role includes everything the one before it can do.

\section{The stages of an incident}

\begin{center}
\begin{tabular}{@{}lll@{}}
\toprule
\textbf{Stage} & \textbf{Means} & \textbf{Can move to} \\
\midrule
Open         & Reported, nobody has responded  & Acknowledged, Resolved \\
Acknowledged & Somebody is working on it       & Mitigated, Resolved \\
Mitigated    & Under control, not fixed        & Resolved, Acknowledged \\
Resolved     & Fixed                           & Closed, Acknowledged \\
Closed       & Finished                        & --- nothing \\
\bottomrule
\end{tabular}
\end{center}

Two rules worth knowing:

\begin{itemize}[leftmargin=1.2cm,itemsep=3pt]
\item Nothing ever returns to \textbf{Open}. Once a person has seen it, the
  ``nobody has looked at this'' state would no longer be true.
\item \textbf{Resolved} can go back to \textbf{Acknowledged} if the problem
  returns. When that happens the recorded resolution time is discarded --- an
  incident that came back was never really resolved.
\end{itemize}

\section{Further reading}

\begin{center}
\begin{tabular}{@{}>{\raggedright\arraybackslash}p{5.4cm}>{\raggedright\arraybackslash}p{9.6cm}@{}}
\toprule
\texttt{README.md} & Overview and quick start \\
\texttt{docs/architecture.md} & Why the system is built this way \\
\texttt{docs/api.md} & Every endpoint, in detail \\
\texttt{docs/openapi.yaml} & The machine-readable API specification \\
\texttt{docs/runbook.md} & Deploying it, and what to do when it breaks \\
\texttt{docs/interview-notes.md} & The design decisions, argued in depth \\
\bottomrule
\end{tabular}
\end{center}

\vfill
\begin{center}
\color{muted}\small
If every checkbox in this manual ticks, the system is working as designed.
\end{center}

\end{document}
